% ===================== 阅读 Reading =====================
\ptesection{阅读}{Reading · FIB + MCM-R + RO + FIBD\&D + MCS-R}{51}

\vspace{1mm}
{\small\color{gray} 本部分共 \textbf{15 题}：\textbf{FIB 读写填空}（下拉选词）Q1--5、\textbf{MCM-R 多选}Q6--7、\textbf{RO 段落重排}Q8--9、\textbf{FIBD\&D 读填空}（拖词入空）Q10--13、\textbf{MCS-R 单选}Q14--15。
每题保留文章原文、你的作答（空处对错与平台一致：\rok{绿勾}=选对，{\color{cWrong}\ding{55}}\,红叉=选错并给出正确答案）、选项/词库（正确项加粗）、你的答案与得分，并为每篇文章附\textbf{中文翻译}。
\textbf{注意（与 2B 不同）}：本套平台"答案\,\&\,评分详情"截图\textbf{未附中文解析}；本卷各题「解析」为\textbf{编者据原文补写}，仅供参考——题干、选项、正确答案、你的作答与得分均忠实誊录自原图。本套阅读官方得分 \textbf{51 / 90}。}
\vspace{2mm}

% ---------------- Q1 ----------------
\begin{qblock}{Q1\quad FIB \#1076\ \ Honorary Fellowship}{读写填空 · 评分 \textcolor{cAccent}{\bfseries 5 / 5}}
\ptesub{文章 Passage（含填空作答）}
\begin{promptbox}
The Royal Geological Society recently \rok{awarded} an honorary fellowship to a pioneering environmental scientist at a prestigious ceremony. Dr. Eliza Montgomery has been recognized for her groundbreaking work in climate change research, \rok{particularly} her studies on the impact of deforestation on global weather patterns. Dr. Montgomery's work has significantly advanced our understanding of how terrestrial ecosystems influence climate dynamics. Her research, which employs both satellite imagery and field data, \rok{has shed} light on the critical role of forests in carbon sequestration and their effect on atmospheric conditions. Beyond her research, Dr. Montgomery has been instrumental in shaping environmental policy through her advisory roles on several international environmental panels. Her \rok{commitment} to environmental education has also led her to mentor the next generation of scientists, emphasizing the importance of sustainable practices. As she accepted her fellowship, Dr. Montgomery stressed the \rok{urgency} of addressing climate change and the responsibility of the scientific community to lead the way in finding viable solutions.
\end{promptbox}
\zh{英国皇家地质学会最近在一场隆重的典礼上，向一位开创性的环境科学家授予了荣誉院士称号。Eliza Montgomery 博士因其在气候变化研究上的开创性工作而受表彰，尤其是她关于森林砍伐对全球天气模式影响的研究。她的工作极大推进了我们对陆地生态系统如何影响气候动态的理解。她的研究同时运用卫星影像与实地数据，揭示了森林在碳封存中的关键作用及其对大气状况的影响。除研究之外，她还通过在多个国际环境专家组中的咨询角色，在制定环境政策方面发挥了重要作用。她对环境教育的投入也促使她指导下一代科学家，强调可持续实践的重要性。在接受院士称号时，Montgomery 博士强调了应对气候变化的紧迫性，以及科学界带头寻找可行方案的责任。}
\ptesub{选项 Choices（每空 4 选 1，正确项加粗）}
{\small
1.\ \textbf{awarded},\ sought,\ validated,\ assessed\\
2.\ \textbf{particularly},\ primarily,\ extremely,\ completely\\
3.\ \textbf{has shed}\quad{\footnotesize\color{gray}（此空及第 4、5 空的四选项在原图中被平台悬浮条遮挡，仅得正确答案）}\\
4.\ \textbf{commitment}\\
5.\ \textbf{urgency}}
\ptesub{解析 Explanation（编者补注）}
\begin{notebox}\footnotesize
\textbf{1.\ awarded}\quad award sth to sb"授予"，主语为学会、宾语为 fellowship，故用 awarded。\par\smallskip
\textbf{2.\ particularly}\quad 引出研究中"尤其"值得一提的具体方向（森林砍伐研究），particularly"尤其"最贴切。\par\smallskip
\textbf{3.\ has shed}\quad 固定搭配 shed light on"揭示、阐明"；主语 Her research 用现在完成时 has shed。\par\smallskip
\textbf{4.\ commitment}\quad commitment to (doing) sth"对……的投入/承诺"，与 to environmental education 搭配。\par\smallskip
\textbf{5.\ urgency}\quad stress the urgency of"强调……的紧迫性"，与后文 addressing climate change 相呼应。
\end{notebox}
\smallskip
\textbf{你的答案：}\ awarded, particularly, has shed, commitment, urgency\qquad\textbf{本题得分：}\ {\color{cAccent}\bfseries 5 / 5}\quad{\footnotesize\color{gray}（用时 02:45，全对）}
\end{qblock}

% ---------------- Q2 ----------------
\begin{qblock}{Q2\quad FIB \#1072\ \ Secret Farm}{读写填空 · 评分 \textcolor{cAccent}{\bfseries 3 / 5}}
\ptesub{文章 Passage（含填空作答）}
\begin{promptbox}
A secret farm in South Africa is \rno{reserve}{home} to 2,000 white rhinos, one of the largest herds in the world. The farm was recently bought by African Parks, a conservation charity that plans to rewild the rhinos and \rno{detach}{restore} them to their natural habitats across Africa. The farm was founded by John Hume, who wanted to \rok{breed} rhinos and sell their horns, but ran out of money. African Parks saw the opportunity to save the rhinos and use them to rebuild populations that have been decimated by poaching. Poaching is driven by the demand for rhino horn in some parts of Asia, where it is used in traditional medicine and as a status symbol. The charity hopes to transfer all the rhinos \rok{from} the farm to secure and sustainable locations in the next 10 years. This would be one of the most ambitious animal rewilding projects ever \rok{undertaken}, and could secure the future of the white rhino in Africa.
\end{promptbox}
\zh{南非一座秘密农场是 2,000 头白犀牛的家园，这是世界上最大的犀牛种群之一。该农场最近被保护慈善机构 African Parks 收购，计划让犀牛重新野化，并把它们放归到非洲各地的自然栖息地。农场由 John Hume 创建，他本想繁育犀牛并出售牛角，却资金耗尽。African Parks 抓住机会拯救这些犀牛，并用它们来重建因偷猎而锐减的种群。偷猎源于亚洲部分地区对犀牛角的需求，那里把它用作传统药材和身份象征。该慈善机构希望在未来 10 年内把所有犀牛从农场转移到安全、可持续的地点。这将是有史以来最雄心勃勃的动物野化项目之一，并可能确保非洲白犀牛的未来。}
\ptesub{选项 Choices（每空 4 选 1，正确项加粗）}
{\small
1.\ ecosystem,\ \textbf{home},\ reserve,\ guardian\\
2.\ detach,\ expel,\ \textbf{restore},\ withdraw\\
3.\ covet,\ \textbf{breed},\ disguise,\ dismantle\\
4.\ \textbf{from},\ away,\ until,\ toward\\
5.\ bypassed,\ \textbf{undertaken},\ overestimated,\ emigrated}
\ptesub{解析 Explanation（编者补注）}
\begin{notebox}\footnotesize
\textbf{1.\ home}\quad be home to"是……的栖息地/家园"，固定表达；reserve"保护区"语义与句式不合。\par\smallskip
\textbf{2.\ restore}\quad restore sb to somewhere"使……回归到"；与 rewild、natural habitats 语义一致。detach"使分离"相反，故你的答案错。\par\smallskip
\textbf{3.\ breed}\quad wanted to breed rhinos and sell their horns"繁育犀牛并卖角"，breed 最贴切。\par\smallskip
\textbf{4.\ from}\quad transfer all the rhinos from the farm to \dots"从农场转移到……"，from\dots to 结构。\par\smallskip
\textbf{5.\ undertaken}\quad projects ever undertaken"曾开展过的项目"，被动完成式修饰 projects。
\end{notebox}
\smallskip
\textbf{你的答案：}\ {\color{cWrong}reserve}, {\color{cWrong}detach}, breed, from, undertaken\qquad\textbf{本题得分：}\ {\color{cAccent}\bfseries 3 / 5}\quad{\footnotesize\color{gray}（用时 02:39；第 1 空 reserve→home、第 2 空 detach→restore）}
\end{qblock}

% ---------------- Q3 ----------------
\begin{qblock}{Q3\quad FIB \#1058\ \ Urban Transformation}{读写填空 · 评分 \textcolor{cAccent}{\bfseries 0 / 4}}
\ptesub{文章 Passage（含填空作答）}
\begin{promptbox}
The focus of this study pivots to the heart of urban transformation, examining the impact of electric scooters in metropolitan areas. In two major cities, a pilot program was \rno{legislated}{introduced} to integrate e-scooters into the urban transportation grid. The study observed how this integration affected traffic flow, pollution levels, and the overall urban landscape. The participating cities saw a marked decrease in car usage, suggesting a shift towards more sustainable travel options. \rno{Totally}{Additionally}, there was a notable reduction in noise and air pollution. The study, named the `Urban Mobility Project', aimed to understand how such small-scale transport means could significantly alter the \rno{architecture}{dynamics} of city life. The results indicated not only a positive environmental impact but also a change in the social fabric of urban communities, \rno{complicating}{fostering} a more connected and accessible city environment.
\end{promptbox}
\zh{本研究的焦点转向城市转型的核心，考察电动滑板车在大都市地区的影响。在两座大城市，一项试点项目被推出，以将电动滑板车纳入城市交通网络。研究观察了这种整合如何影响交通流量、污染水平和整体城市面貌。参与城市的汽车使用明显减少，显示出向更可持续出行方式的转变。此外，噪音和空气污染也明显下降。这项名为"城市出行项目"的研究，旨在了解这类小型交通工具如何能显著改变城市生活的运作格局。结果表明，它不仅带来积极的环境影响，也改变了城市社区的社会结构，促进了一个更互联、更便捷的城市环境。}
\ptesub{选项 Choices（每空 4 选 1，正确项加粗）}
{\small
1.\ legislated,\ sifted,\ participated,\ \textbf{introduced}\\
2.\ Nevertheless,\ Totally,\ \textbf{Additionally},\ Firstly\\
3.\ architecture,\ \textbf{dynamics},\ boundaries,\ engagement\\
4.\ surpassing,\ complicating,\ \textbf{fostering},\ admitting}
\ptesub{解析 Explanation（编者补注）}
\begin{notebox}\footnotesize
\textbf{1.\ introduced}\quad a pilot program was introduced"推出试点项目"；legislated"立法"与 program 不搭。\par\smallskip
\textbf{2.\ Additionally}\quad 上句讲汽车使用减少，本句再添一项好处（噪音与污染下降），用 Additionally"此外"递进；Totally"完全地"不合逻辑。\par\smallskip
\textbf{3.\ dynamics}\quad alter the dynamics of city life"改变城市生活的运作格局/动态"；architecture"建筑结构"过于具体，不合语境。\par\smallskip
\textbf{4.\ fostering}\quad fostering a more connected environment"促进更互联的环境"，褒义；complicating"使复杂化"贬义，与全句积极基调矛盾，故你的答案错。
\end{notebox}
\smallskip
\textbf{你的答案：}\ {\color{cWrong}legislated, Totally, architecture, complicating}\qquad\textbf{本题得分：}\ {\color{cAccent}\bfseries 0 / 4}\quad{\footnotesize\color{gray}（用时 02:05；四空全错，正确为 introduced / Additionally / dynamics / fostering）}
\end{qblock}

% ---------------- Q4 ----------------
\begin{qblock}{Q4\quad FIB \#1045\ \ Social Media}{读写填空 · 评分 \textcolor{cAccent}{\bfseries 3 / 4}}
\ptesub{文章 Passage（含填空作答）}
\begin{promptbox}
Social media's impact on the youth has been a topic of considerable debate. Platforms like Instagram, Twitter, and Facebook offer spaces for social interaction, self-expression, and information sharing. \rok{However}, they also pose risks such as cyberbullying, exposure to \rok{inappropriate} content, and the impact on mental health. Studies have shown that excessive use of social media can lead to feelings of anxiety, depression, and loneliness in young people. The pressure to maintain a certain image online and the constant comparison with others can exacerbate these feelings. On the positive side, social media can also \rno{subscribe}{provide} support networks, educational content, and opportunities for civic engagement among youth. The challenge lies in achieving a balance, \rok{promoting} safe and responsible use of these platforms.
\end{promptbox}
\zh{社交媒体对青少年的影响一直是备受争议的话题。Instagram、Twitter、Facebook 等平台提供了社交互动、自我表达和信息分享的空间。然而，它们也带来风险，如网络霸凌、接触不当内容以及对心理健康的影响。研究表明，过度使用社交媒体会让年轻人产生焦虑、抑郁和孤独感。在网上维持某种形象的压力以及不断与他人比较，会加剧这些情绪。从积极方面看，社交媒体也能为青少年提供支持网络、教育内容和参与公共事务的机会。挑战在于取得平衡，促进对这些平台安全、负责任的使用。}
\ptesub{选项 Choices（每空 4 选 1，正确项加粗）}
{\small
1.\ Therefore,\ That is,\ Similarly,\ \textbf{However}\\
2.\ ridiculous,\ humorous,\ \textbf{inappropriate},\ underrated\\
3.\ \textbf{provide},\ trigger,\ validate,\ subscribe\\
4.\ delegating,\ experimenting,\ \textbf{promoting},\ voting}
\ptesub{解析 Explanation（编者补注）}
\begin{notebox}\footnotesize
\textbf{1.\ However}\quad 前句讲平台的好处、本句讲风险，语义转折，用 However"然而"。\par\smallskip
\textbf{2.\ inappropriate}\quad exposure to inappropriate content"接触不当内容"，与 cyberbullying 等风险并列。\par\smallskip
\textbf{3.\ provide}\quad social media can provide support networks"提供支持网络"；subscribe"订阅"不能带此类宾语，故你的答案错。\par\smallskip
\textbf{4.\ promoting}\quad promoting safe and responsible use"促进安全负责任的使用"，与 achieving a balance 呼应。
\end{notebox}
\smallskip
\textbf{你的答案：}\ However, inappropriate, {\color{cWrong}subscribe}, promoting\qquad\textbf{本题得分：}\ {\color{cAccent}\bfseries 3 / 4}\quad{\footnotesize\color{gray}（用时 00:53；第 3 空 subscribe→provide）}
\end{qblock}

% ---------------- Q5 ----------------
\begin{qblock}{Q5\quad FIB \#1040\ \ Drought-resistant Crops}{读写填空 · 评分 \textcolor{cAccent}{\bfseries 1 / 4}}
\ptesub{文章 Passage（含填空作答）}
\begin{promptbox}
The impact of drought-resistant crops could be transformative, particularly for smallholder farmers in developing countries, who are often the most vulnerable to the effects of climate change. By \rno{resurrecting}{cultivating} these crops, farmers can sustain their livelihoods and provide food for their communities even in challenging conditions. However, the adoption of these crops also raises concerns. Issues such as biodiversity \rno{conflict}{loss}, dependency on specific seed companies, and the long-term ecological impacts of genetically modified organisms are part of an ongoing debate. Despite these \rok{concerns}, the potential benefits of drought-resistant crops are significant. They represent a crucial step in adapting our agricultural systems to the changing climate. As research continues and these crops become more sophisticated, they could play a key role in ensuring global food \rno{emergence}{security} in an increasingly unpredictable environment.
\end{promptbox}
\zh{抗旱作物的影响可能是变革性的，尤其对发展中国家的小农户而言——他们往往最易受气候变化影响。通过种植这些作物，农民即便在艰难条件下也能维持生计、为社区提供食物。然而，采用这些作物也引发了担忧。诸如生物多样性丧失、对特定种子公司的依赖，以及转基因生物的长期生态影响等问题，仍是持续争论的一部分。尽管存在这些担忧，抗旱作物的潜在好处依然显著。它们是让农业系统适应气候变化的关键一步。随着研究推进、这些作物日益成熟，它们在日益难以预测的环境中，可能对确保全球粮食安全发挥关键作用。}
\ptesub{选项 Choices（每空 4 选 1，正确项加粗）}
{\small
1.\ resurrecting,\ demonstrating,\ \textbf{cultivating},\ devastating\\
2.\ \textbf{loss},\ conflict,\ shift,\ draft\\
3.\ \textbf{concerns},\ facilities,\ editions,\ fees\\
4.\ \textbf{security},\ inevitability,\ support,\ emergence}
\ptesub{解析 Explanation（编者补注）}
\begin{notebox}\footnotesize
\textbf{1.\ cultivating}\quad By cultivating these crops"通过种植这些作物"；resurrecting"使复活"与 crops 不搭，故你的答案错。\par\smallskip
\textbf{2.\ loss}\quad biodiversity loss"生物多样性丧失"，是常见的环境隐患表达；conflict"冲突"不合，故你的答案错。\par\smallskip
\textbf{3.\ concerns}\quad Despite these concerns"尽管有这些担忧"，回指前文 raises concerns。\par\smallskip
\textbf{4.\ security}\quad global food security"全球粮食安全"固定搭配；emergence"出现"不合，故你的答案错。
\end{notebox}
\smallskip
\textbf{你的答案：}\ {\color{cWrong}resurrecting, conflict}, concerns, {\color{cWrong}emergence}\qquad\textbf{本题得分：}\ {\color{cAccent}\bfseries 1 / 4}\quad{\footnotesize\color{gray}（用时 01:28；仅第 3 空 concerns 正确）}
\end{qblock}

% ---------------- Q6 ----------------
\begin{qblock}{Q6\quad MCM-R \#23}{多选 · 评分 \textcolor{cAccent}{\bfseries 0 / 2}}
{\footnotesize\color{gray} 注：原图标题处仅重复显示题型"MCM-R"，无独立标题；下文按内容暂记为"Psychodynamic Approach"。}
\ptesub{文章 Passage}
\begin{promptbox}
The Psychodynamic Approach. Theorists adopting the psychodynamic approach hold that inner conflicts are crucial for understanding human behavior, including aggression. Sigmund Freud, for example, believed that aggressive impulses are inevitable reactions to the frustrations of daily life. Children normally desire to vent aggressive impulses on other people, including their parents, because even the most attentive parents cannot gratify all of their demands immediately. Yet children, also fearing their parents' punishment and the loss of parental love, come to repress most aggressive impulses. The Freudian perspective, in a sense: sees us as ``steam engines.'' By holding in rather than venting ``steam,'' we set the stage for future explosions. Pent-up aggressive impulses demand outlets. They may be expressed toward parents in indirect ways such as destroying furniture, or they may be expressed toward strangers later in life.
\end{promptbox}
\zh{心理动力学取向。采用心理动力学取向的理论家认为，内在冲突对于理解人类行为（包括攻击性）至关重要。例如，弗洛伊德认为，攻击性冲动是对日常生活挫折不可避免的反应。儿童通常渴望把攻击性冲动发泄到他人身上，包括父母，因为即使最细心的父母也无法立即满足他们的一切要求。然而，儿童同时害怕父母的惩罚以及失去父母之爱，于是逐渐压抑大多数攻击性冲动。从某种意义上说，弗洛伊德的观点把我们看作"蒸汽机"：通过憋住而非释放"蒸汽"，我们为将来的爆发埋下伏笔。被压抑的攻击性冲动需要出口，它们可能以间接方式发泄到父母身上（如破坏家具），也可能在日后发泄到陌生人身上。}
\ptesub{题目 Question}
\begin{notebox}According to the paragraph, Freud believed that children experience conflict between a desire to vent aggression on their parents and\end{notebox}
\ptesub{选项 Choices（正确项绿色加粗）}
{\small
A) a frustration that their parents do not give them everything they want\\
{\color{cFix}\textbf{B) a fear that their parents will punish them}}\\
C) a desire to take care of their parents\\
D) a desire to vent aggression on other family members\\
{\color{cFix}\textbf{E) a fear for loss of their parents' love}}}
\par\smallskip
\textbf{正确答案}\ \textcolor{cFix}{\bfseries B, E}\qquad\textbf{你的答案}\ {\color{cWrong}\bfseries C}\qquad\textbf{本题得分：}\ {\color{cAccent}\bfseries 0 / 2}\quad{\footnotesize\color{gray}（用时 00:09）}
\end{qblock}

% ---------------- Q7 ----------------
\begin{qblock}{Q7\quad MCM-R \#34\ \ Paintings}{多选 · 评分 \textcolor{cAccent}{\bfseries 0 / 2}}
\ptesub{文章 Passage}
\begin{promptbox}
A third opinion takes psychological motivation much further into the realm of tribal ceremonies and mystery: the belief that certain animals assumed mythical significance as ancient ancestors or protectors of a given tribe or clan. Two types of images substantiate this theory: the strange, indecipherable geometric shapes that appear near some animals, and the few drawings of men. Wherever men appear they are crudely drawn and their bodies are elongated and rigid. Some men are in a prone position and some have bird or animal heads. Advocates for this opinion point to reports from people who have experienced a trance state, a highly suggestive state of low consciousness between waking and sleeping. Uniformly, these people experienced weightlessness and the sensation that their bodies were being stretched lengthwise. Advocates also point to people who believe that the forces of nature are inhabited by spirits, particularly shamans who believe that an animal's spirit and energy is transferred to them while in a trance. One Lascaux narrative picture, which shows a man with a birdlike head and a wounded animal, would seem to lend credence to this third opinion, but there is still much that remains unexplained. For example, where is the proof that the man in the picture is a shaman? He could as easily be a hunter wearing a headmask. Many tribal hunters, including some Native Americans, camouflaged themselves by wearing animal heads and hides.
\end{promptbox}
\zh{第三种观点把心理动机进一步引向部落仪式与神秘领域：即认为某些动物作为特定部落或氏族的远古祖先或守护者，具有神话意义。两类图像佐证了这一理论：出现在某些动物旁边的奇异、难以解读的几何图形，以及少数人物画。凡有人物出现，都画得粗糙，身体被拉长且僵硬；有些人呈俯卧姿势，有些则长着鸟或动物的头。该观点的支持者援引了经历过"入迷状态"者的描述——那是一种介于清醒与睡眠之间、高度受暗示、意识低下的状态。这些人无一例外地体验到失重感，以及身体被纵向拉长的感觉。支持者还援引那些相信自然力量由精灵栖居的人，尤其是萨满——他们相信在入迷状态下，动物的灵魂与能量会转移到自己身上。拉斯科的一幅叙事性壁画描绘了一个长着鸟头的人和一头受伤的动物，似乎为第三种观点提供了佐证，但仍有许多未解之处。例如，凭什么证明画中人是萨满？他也完全可能是戴着头面具的猎人。许多部落猎人，包括一些北美原住民，都会披戴兽头兽皮来伪装自己。}
\ptesub{题目 Question}
\begin{notebox}According to the paragraph, why do some scholars refer to a trance state to help understand the cave paintings?\end{notebox}
\ptesub{选项 Choices（正确项绿色加粗）}
{\small
A) To explain the state of consciousness the artists were in when they painted their pictures\\
B) To demonstrate the mythical significance of the strange geometric shapes\\
C) To indicate that trance states were often associated with activities that took place inside caves\\
{\color{cFix}\textbf{D) To give a possible reason for the strange appearance of the men painted on the cave walls}}\\
{\color{cFix}\textbf{E) To substantiate the belief that certain animals take mythical significance.}}}
\par\smallskip
\textbf{正确答案}\ \textcolor{cFix}{\bfseries D, E}\qquad\textbf{你的答案}\ {\color{cWrong}\bfseries B}\qquad\textbf{本题得分：}\ {\color{cAccent}\bfseries 0 / 2}\quad{\footnotesize\color{gray}（用时 00:04）}
\end{qblock}

% ---------------- Q8 ----------------
\begin{qblock}{Q8\quad RO \#835\ \ Livestock}{段落重排 · 评分 \textcolor{cAccent}{\bfseries 3 / 3}}
\ptesub{乱序段落 Paragraphs}
\begin{promptbox}
\textbf{1)} However, intensive livestock farming raises concerns about environmental sustainability and animal welfare.\\
\textbf{2)} Thus, there is a growing emphasis on sustainable farming practices.\\
\textbf{3)} Livestock play a crucial role in agriculture.\\
\textbf{4)} They not only provide meat and dairy products but also contribute to crop production through manure, which enriches soil fertility.
\end{promptbox}
\zh{（按正确顺序）牲畜在农业中起着关键作用；它们不仅提供肉类和乳制品，还通过粪肥促进作物生产、增加土壤肥力；然而，集约化畜牧引发了对环境可持续性和动物福利的担忧；因此，人们越来越强调可持续的养殖方式。}
\smallskip
\textbf{正确顺序}\ \textcolor{cFix}{\bfseries 3 → 4 → 1 → 2}\qquad\textbf{你的顺序}\ \textcolor{cFix}{\bfseries 3 → 4 → 1 → 2}\qquad\textbf{本题得分：}\ {\color{cAccent}\bfseries 3 / 3}\quad{\footnotesize\color{gray}（用时 00:44，全对）}
\end{qblock}

% ---------------- Q9 ----------------
\begin{qblock}{Q9\quad RO \#827\ \ Fossil Fuel}{段落重排 · 评分 \textcolor{cAccent}{\bfseries 3 / 3}}
\ptesub{乱序段落 Paragraphs}
\begin{promptbox}
\textbf{1)} Thus, the resulting climate change concerns have catalyzed the search for cleaner, more sustainable energy alternatives.\\
\textbf{2)} Their abundant energy and ease of extraction and transport made them the preferred choice.\\
\textbf{3)} Fossil fuels, including coal, oil, and natural gas, have been the cornerstone of industrial development for over a century.\\
\textbf{4)} However, their combustion releases significant greenhouse gases, leading to global warming and environmental degradation.
\end{promptbox}
\zh{（按正确顺序）化石燃料——包括煤、石油和天然气——一个多世纪以来一直是工业发展的基石；它们能量丰富、易于开采和运输，因而成为首选；然而，其燃烧会释放大量温室气体，导致全球变暖和环境退化；因此，由此产生的气候变化担忧，催生了对更清洁、更可持续的替代能源的探索。}
\smallskip
\textbf{正确顺序}\ \textcolor{cFix}{\bfseries 3 → 2 → 4 → 1}\qquad\textbf{你的顺序}\ \textcolor{cFix}{\bfseries 3 → 2 → 4 → 1}\qquad\textbf{本题得分：}\ {\color{cAccent}\bfseries 3 / 3}\quad{\footnotesize\color{gray}（用时 00:14，全对）}
\end{qblock}

% ---------------- Q10 ----------------
\begin{qblock}{Q10\quad FIBD\&D \#1110\ \ Deep-sea Fish}{读填空（拖词）· 评分 \textcolor{cAccent}{\bfseries 4 / 4}}
\ptesub{文章 Passage（含填空作答）}
\begin{promptbox}
The world of science has witnessed countless breakthroughs. Last month, a research team \rok{discovered} a new species of deep-sea fish during an expedition. These creatures, living in the darkest depths, were hidden from human sight until now. The scientists didn't just \rok{stumble} across them casually; they used advanced sonar technology and spent hours analyzing data. Back in the lab, the team worked with biologists who \rok{possessed} extensive knowledge of marine genetics. Their combined expertise enabled them to study the fish's unique DNA structure. Soon after, their findings were published, and the scientific community \rok{hailed} it a significant discovery, opening new doors for ocean research.
\end{promptbox}
\zh{科学界见证了无数突破。上个月，一支研究团队在一次考察中发现了一种新的深海鱼。这些生物生活在最黑暗的深处，此前一直不为人所见。科学家并非只是随意地偶然发现它们；他们使用了先进的声呐技术，并花了数小时分析数据。回到实验室后，团队与拥有丰富海洋遗传学知识的生物学家合作。他们综合的专业知识使其得以研究这种鱼独特的 DNA 结构。不久之后，他们的发现得以发表，科学界称赞这是一项重大发现，为海洋研究开启了新的大门。}
\ptesub{词库 Word Bank（拖词入空，正确答案加粗）}
{\small 1) \textbf{discovered}\quad 2) \textbf{possessed}\quad 3) invented\quad 4) find\quad 5) \textbf{hailed}\quad 6) \textbf{stumble}\quad 7) occupied}
\ptesub{解析 Explanation（编者补注）}
\begin{notebox}\footnotesize
\textbf{1.\ discovered}\quad a research team discovered a new species"发现新物种"，一般过去时。\par\smallskip
\textbf{2.\ stumble}\quad didn't just stumble across them"并非只是偶然发现"；助动词 did 之后用原形 stumble。\par\smallskip
\textbf{3.\ possessed}\quad biologists who possessed extensive knowledge"拥有丰富知识的生物学家"，定语从句谓语。\par\smallskip
\textbf{4.\ hailed}\quad hailed it a significant discovery"称赞它是重大发现"。（干扰词 invented / find / occupied 均不合。）
\end{notebox}
\smallskip
\textbf{你的答案：}\ discovered, stumble, possessed, hailed\qquad\textbf{本题得分：}\ {\color{cAccent}\bfseries 4 / 4}\quad{\footnotesize\color{gray}（用时 01:44，全对）}
\end{qblock}

% ---------------- Q11 ----------------
\begin{qblock}{Q11\quad FIBD\&D \#1109\ \ Music Show}{读填空（拖词）· 评分 \textcolor{cAccent}{\bfseries 4 / 4}}
\ptesub{文章 Passage（含填空作答）}
\begin{promptbox}
Last night, I went to a grand music show at the city hall. The auditorium was filled with excited audiences. Some people, unfortunately, were chatting \rok{constantly}, ignoring the performers onstage, seemingly unaware of the rudeness. \rok{However}, I was determined not to let that spoil my mood. I closed my eyes, \rok{trying} to block out the noise, and still tried my best to focus on the soulful music, \rok{immersing} myself in the wonderful melodies.
\end{promptbox}
\zh{昨晚，我去市政厅看了一场盛大的音乐演出。礼堂里坐满了兴奋的观众。不幸的是，有些人一直在聊天，无视台上的表演者，似乎意识不到这样很失礼。然而，我下定决心不让这破坏我的心情。我闭上眼睛，努力屏蔽噪音，仍尽力专注于那充满感情的音乐，让自己沉浸在美妙的旋律之中。}
\ptesub{词库 Word Bank（拖词入空，正确答案加粗）}
{\small 1) \textbf{immersing}\quad 2) \textbf{constantly}\quad 3) Therefore\quad 4) \textbf{However}\quad 5) immersed\quad 6) \textbf{trying}\quad 7) tries\quad 8) insignificantly}
\ptesub{解析 Explanation（编者补注）}
\begin{notebox}\footnotesize
\textbf{1.\ constantly}\quad were chatting constantly"一直在聊天"，副词修饰进行时。\par\smallskip
\textbf{2.\ However}\quad 上句讲他人无礼、本句讲"我"不受影响，语义转折用 However；Therefore"因此"不合。\par\smallskip
\textbf{3.\ trying}\quad I closed my eyes, trying to block out the noise"闭眼、努力屏蔽噪音"，现在分词作伴随状语。\par\smallskip
\textbf{4.\ immersing}\quad immersing myself in the melodies"沉浸于旋律"，现在分词伴随；immersed 为过去式/过去分词，此处需分词形式。
\end{notebox}
\smallskip
\textbf{你的答案：}\ constantly, However, trying, immersing\qquad\textbf{本题得分：}\ {\color{cAccent}\bfseries 4 / 4}\quad{\footnotesize\color{gray}（用时 00:30，全对）}
\end{qblock}

% ---------------- Q12 ----------------
\begin{qblock}{Q12\quad FIBD\&D \#1108\ \ Sufficient Sleep}{读填空（拖词）· 评分 \textcolor{cAccent}{\bfseries 3 / 4}}
\ptesub{文章 Passage（含填空作答）}
\begin{promptbox}
In our fast-paced modern lives, getting sufficient sleep is often overlooked, yet it \rok{is} vital. And also, it is really important that you get enough sleep to stay \rno{rejuvenated}{rested}. Adequate sleep allows our body to repair and rejuvenate, keeping our mind sharp. If you \rok{lack} enough sleep regularly, you may feel drowsy, which \rok{impacts} concentration, saps energy, and even dulls reflexes, making daily tasks a struggle.
\end{promptbox}
\zh{在快节奏的现代生活中，充足的睡眠常被忽视，但它至关重要。而且，你获得足够睡眠以保持精神饱满，这确实很重要。充足的睡眠让身体得以修复和恢复，使头脑保持敏锐。如果你经常睡眠不足，就可能感到困倦，这会影响注意力、消耗精力，甚至使反应迟钝，让日常任务变得吃力。}
\ptesub{词库 Word Bank（拖词入空，正确答案加粗）}
{\small 1) constitute\quad 2) \textbf{is}\quad 3) rejuvenated\quad 4) \textbf{lack}\quad 5) \textbf{rested}\quad 6) shun\quad 7) \textbf{impacts}}
\ptesub{解析 Explanation（编者补注）}
\begin{notebox}\footnotesize
\textbf{1.\ is}\quad yet it \_\_\_ vital"但它至关重要"，主语 it、系动词 is。\par\smallskip
\textbf{2.\ rested}\quad stay rested"保持得到休息/精神饱满"；stay 为半系动词接形容词，rested 由动词转化而来更贴合语境；rejuvenated 语义偏"返老还童/焕发"，此处平台判为错，故你的答案错。\par\smallskip
\textbf{3.\ lack}\quad If you lack enough sleep"若缺乏足够睡眠"，主语 you、谓语原形 lack。\par\smallskip
\textbf{4.\ impacts}\quad which impacts concentration"影响注意力"，which 指代前句、谓语三单 impacts。
\end{notebox}
\smallskip
\textbf{你的答案：}\ is, {\color{cWrong}rejuvenated}, lack, impacts\qquad\textbf{本题得分：}\ {\color{cAccent}\bfseries 3 / 4}\quad{\footnotesize\color{gray}（用时 01:08；第 2 空 rejuvenated→rested）}
\end{qblock}

% ---------------- Q13 ----------------
\begin{qblock}{Q13\quad FIBD\&D \#1068\ \ Chemical Cauldron}{读填空（拖词）· 评分 \textcolor{cAccent}{\bfseries 2 / 4}}
\ptesub{文章 Passage（含填空作答）}
\begin{promptbox}
The universe is a vast chemical cauldron, brewing elements essential for life. Among the 92 naturally occurring elements, a select few form the building blocks of living \rok{organisms}. Oxygen, carbon, hydrogen, and nitrogen, for instance, constitute the \rok{majority} of life's molecular structure. This universality of life's elemental ingredients suggests that the search for extraterrestrial life, \rno{despite}{regardless} of its chemical basis, may find success in the far \rno{prospect}{reaches} of our galaxy and beyond.
\end{promptbox}
\zh{宇宙是一口巨大的化学熔炉，酝酿着生命所必需的元素。在 92 种天然存在的元素中，少数几种构成了生命有机体的基本组成单元。例如，氧、碳、氢、氮就构成了生命分子结构的大部分。生命元素成分的这种普遍性表明，寻找地外生命——无论其化学基础如何——都可能在我们星系乃至更远处的深远之地取得成功。}
\ptesub{词库 Word Bank（拖词入空，正确答案加粗）}
{\small 1) prospect\quad 2) \textbf{regardless}\quad 3) departments\quad 4) supplements\quad 5) \textbf{organisms}\quad 6) \textbf{reaches}\quad 7) \textbf{majority}\quad 8) despite}
\ptesub{解析 Explanation（编者补注）}
\begin{notebox}\footnotesize
\textbf{1.\ organisms}\quad building blocks of living organisms"生命有机体的基本单元"。\par\smallskip
\textbf{2.\ majority}\quad constitute the majority of\dots"构成……的大部分"。\par\smallskip
\textbf{3.\ regardless}\quad regardless of its chemical basis"无论其化学基础如何"，regardless of 为固定搭配；despite 后不接 of，故你的答案错。\par\smallskip
\textbf{4.\ reaches}\quad the far reaches of our galaxy"星系的深远之处"，reaches 作名词指"辽远的地带"；prospect"前景"与 far\dots of 结构不合，故你的答案错。
\end{notebox}
\smallskip
\textbf{你的答案：}\ organisms, majority, {\color{cWrong}despite, prospect}\qquad\textbf{本题得分：}\ {\color{cAccent}\bfseries 2 / 4}\quad{\footnotesize\color{gray}（用时 02:08；第 3 空 despite→regardless、第 4 空 prospect→reaches）}
\end{qblock}

% ---------------- Q14 ----------------
\begin{qblock}{Q14\quad MCS-R \#45}{单选 · 评分 \textcolor{cAccent}{\bfseries 1 / 1}}
{\footnotesize\color{gray} 注：原图标题处仅重复显示题型"MCS-R"，无独立标题；下文按内容暂记为"Native American Renaissance"。}
\ptesub{文章 Passage}
\begin{promptbox}
The outpouring of contemporary American Indian literature in the last two decades, often called the Native American Renaissance, represents for many the first opportunity to experience Native American poetry. The appreciation of traditional oral American Indian literature has been limited, hampered by poor translations and by the difficulty, even in the rare culturally sensitive and aesthetically satisfying translation, of completely conveying the original's verse structure, tone, and syntax. By writing in English and experimenting with European literary forms, contemporary American Indian writers have broadened their potential audience, while clearly retaining many essential characteristics of their ancestral oral traditions. For example, Pulitzer-prizewinning author N. Scott Momaday's poetry often treats art and mortality in a manner that recalls British romantic poetry, while his poetic response to the power of natural forces recalls Cherokee oral literature. In the same way, his novels, an art form European in origin, display an eloquence that echoes the oratorical grandeur of the great nineteenth-century American Indian chiefs.
\end{promptbox}
\zh{近二十年来当代美洲印第安文学的大量涌现，常被称为"美洲原住民文艺复兴"，对许多人而言，这是首次体验美洲原住民诗歌的机会。人们对传统口传印第安文学的欣赏一直有限，既受制于拙劣的翻译，也因为即便是罕见的、兼具文化敏感度与美学感染力的译本，也难以完整传达原作的诗律结构、语气与句法。通过用英语写作并尝试欧洲文学形式，当代印第安作家拓宽了潜在读者群，同时又清晰地保留了祖先口传传统的诸多本质特征。例如，普利策奖得主 N. Scott Momaday 的诗歌常以令人想起英国浪漫主义诗歌的方式处理艺术与死亡，而他对自然力量的诗意回应则让人想起切罗基口传文学。同样，他那源自欧洲艺术形式的小说，展现出一种雄辩之力，呼应着 19 世纪伟大的美洲印第安酋长们的演说气魄。}
\ptesub{题目 Question}
\begin{notebox}Which of the following is most likely one of the reasons that the author mentions the work of N. Scott Momaday?\end{notebox}
\ptesub{选项 Choices}
{\small
A)\ {\color{gray}（原图此处被平台悬浮条遮挡，选项文字未能获取；此项即正确答案）}\\
B) To emphasize the similarities between Momaday's writings and their European literary models\\
C) To demonstrate the contemporary appeal of traditional Native American oral literature\\
D) To suggest that contemporary American Indian writers have sacrificed traditional values for popular literary success}
\par\smallskip
\textbf{正确答案}\ \textcolor{cFix}{\bfseries A}\ {\footnotesize\color{gray}（文字被遮挡，见上注）}\qquad\textbf{你的答案}\ \textcolor{cFix}{\bfseries A}\qquad\textbf{本题得分：}\ {\color{cAccent}\bfseries 1 / 1}\quad{\footnotesize\color{gray}（用时 00:03，正确）}
\end{qblock}

% ---------------- Q15 ----------------
\begin{qblock}{Q15\quad MCS-R \#38}{单选 · 评分 \textcolor{cAccent}{\bfseries 0 / 1}}
{\footnotesize\color{gray} 注：原图标题处仅重复显示题型"MCS-R"，无独立标题；下文按内容暂记为"Pinnipeds' Flippers"。}
\ptesub{文章 Passage}
\begin{promptbox}
Biologists have long maintained that two groups of pinnipeds, sea lions and walruses, are descended from a terrestrial bearlike animal, whereas the remaining group, seals, shares an ancestor with weasels. But the recent discovery of detailed similarities in the skeletal structure of the flippers in all three groups undermines the attempt to explain away superficial resemblance as due to convergent evolution---the independent development of similarities between unrelated groups in response to similar environmental pressures. Flippers may indeed be a necessary response to aquatic life; turtles, whales, and dugongs also have them. But the common detailed design found among the pinnipeds probably indicates a common ancestor. Moreover, walruses and seals drive themselves through the water with thrusts of their hind flippers, but sea lions use their front flippers. If anatomical similarity in the flippers resulted from similar environmental pressures, as posited by the convergent evolution theory, one would expect walruses and seals, but not seals and sea lions, to have similar flippers.
\end{promptbox}
\zh{长期以来，生物学家认为两类鳍足动物——海狮和海象——是从一种陆生的熊状动物演化而来，而其余一类海豹则与黄鼠狼共有一个祖先。但最近发现，三类动物鳍肢骨骼结构存在细致的相似之处，这动摇了把表面相似解释为"趋同进化"（即无亲缘关系的类群为应对相似的环境压力而各自独立发展出相似特征）的企图。鳍肢或许确实是对水生生活的必要适应；海龟、鲸和儒艮也都有鳍。但鳍足动物之间共有的这种细致构造，很可能表明它们有共同祖先。此外，海象和海豹靠后鳍推水前进，而海狮则用前鳍。若鳍肢的解剖学相似确如趋同进化理论所设想的那样源于相似的环境压力，那么人们应当预期海象与海豹（而非海豹与海狮）拥有相似的鳍肢。}
\ptesub{题目 Question}
\begin{notebox}According to the passage, it has been recently discovered that\end{notebox}
\ptesub{选项 Choices（正确项绿色加粗）}
{\small
{\color{cFix}\textbf{A) there are detailed skeletal similarities in the flippers of pinnipeds}}\\
B) sea lions, seals, and walruses are all pinnipeds\\
C) pinnipeds are descended from animals that once lived on land\\
D) animals without common ancestors sometimes evolve in similar ways}
\par\smallskip
\textbf{正确答案}\ \textcolor{cFix}{\bfseries A}\qquad\textbf{你的答案}\ {\color{cWrong}\bfseries B}\qquad\textbf{本题得分：}\ {\color{cAccent}\bfseries 0 / 1}\quad{\footnotesize\color{gray}（用时 00:03）}
\end{qblock}
